\section{A uniform tail for spatially integrated stopped densities}
\label{mesh:section}

We prove the estimate that determines the logarithmic correction in the
gauge.  The essential feature is uniformity in the final mesh: a large
threshold is tested first at an intermediate mesh, and the remaining
increments must be controlled down to any prescribed finer mesh.

Throughout this section, $0<\kappa<8$ is fixed,
\[
 p=1-\kappa/8,\qquad d=2-p,\qquad a_0=8/\kappa-1.
\]
We use the capacity normalization
$\partial_tg_t=2/(g_t-\sqrt\kappa B_t)$.  The initial domains considered
here are surviving domains of an ordinary half-plane trace at finite
stopping times, and their Euclidean similarities.  Conditional statements below refer
to the ordinary chordal law in that domain, with its marked tip and
target.  We write $D_t$ for its subsequent surviving domain.

\paragraph{Inputs used here.}
Lemmas~\ref{ll:finite-tilt} and~\ref{ll:approach} supply the nonnegative density
\[
 M_t(z)=R_{D_t}(z)^{-p}\sin^{a_0}\theta_t(z),
\]
continued by zero on killing.  Here $R$ is the full conformal radius and
$\theta$ is the argument of the image point in tip--target half-plane
coordinates.  We use the following precise parts of that law.
First, on a region where the distance from $z$ to the domain boundary is
bounded below, $M_t(z)$ is a bounded stopped martingale.  If a smaller
distance is never reached, its terminal contribution is zero.  Second,
at a live stopping time, the probability of a subsequent approach to
distance $u$, for $u$ smaller than a fixed multiple of the current
distance to the boundary, is at most $C u^pM_t(z)$.  The expected density
at that approach is $M_t(z)$.  All these assertions use only finite
stopped changes of law and monotone exhaustion.  The coefficient of
$dB_t$ in $M_t(z)$ has absolute value at most
\begin{equation}
 C R_{D_t}(z)^{-p}\sin^{a_0}\theta_t(z)/|Z_t(z)|,
 \qquad Z_t(z)=G_t(z)-\sqrt\kappa B_t.
 \label{mesh:coefficient-input}
\end{equation}
Here $G_t$ includes the initial uniformization: if the initial domain
is $J(\mathbb H)$, then $G_t=g_t\circ J^{-1}$, with $g_t$ in the
half-plane Loewner coordinates of Section~\ref{ll:section}.
When $\kappa>4$, Proposition~\ref{ll:disk-rule} supplies one more
estimate.  Suppose $\overline B(z,r)$ is contained in the initial
domain.  At its first circle hit, if $z$ is still live, the conditional
probability that the next exit of $B(z,Lr)$ occurs with $z$ still live
is at most $CL^{-c}$ for all sufficiently large $L$.  The outer disk need
not be free.  The domain class is exactly the one stated above.

After a translation, fix $A=B(0,1)$ with a fixed enlargement contained in
the initial domain.  All the points in $A$ consequently have a uniform
initial boundary clearance, which we denote by $\rho_{\rm col}>0$.
For $0<r<\rho_{\rm col}$, let
$\tau_r(z)$ be the first live approach of the trace to distance $r$
from $z$, with value infinity if there is none, and set
\[
 W_r(z)=M_{\tau_r(z)}(z)\mathbf1_{\{\tau_r(z)<\infty\}},
 \qquad I_r(z)=\mathbf1_{\{\tau_r(z)<\infty\}}.
\]
The maps $(z,\omega)\mapsto\tau_r(z),W_r(z),I_r(z)$ have jointly
measurable versions: before killing this follows from continuous trace
distance and the continuous Loewner maps. The weights $W_r$ and
indicators $I_r$ are zero when no live approach occurs; the stopping
time then has value $\infty$.  One may first localize the maps on
compact subsets of the surviving domain and then exhaust.
On a finite approach, \eqref{ll:trace-clearance} gives
\[
 \operatorname{dist}(z,\partial D_{\tau_r(z)})=r.
\]
Thus at $\tau_r(z)$ the restarted estimates of Lemma~\ref{ll:approach}
apply to the finer distance $qr$ whenever $0<q\le1/(8\mathrm e)$.

\begin{proposition}[Integrated mesh tail]\label{mesh:tail}
There are constants $C,c>0$ and $0<\delta_0<\rho_{\rm col}$, depending only on $\kappa$ and
the fixed collar, such that, for $0<r_f\le\delta\le\delta_0$,
\begin{equation}
 \mathbb P\left\{\int_A W_{r_f}(z)\,dA(z)>C\delta^{-p}\right\}
 \le C\exp(-c\delta^{-2}).
 \label{mesh:tail-equation}
\end{equation}
The constants are uniform over the stated initial marked domains, including
conditioning on any past for which the collar is free.  They do not
depend on the final distance $r_f$ or on an upper bound for the
distance to the tip.  The same estimate holds if only approaches by a
fixed or stopping time are counted.
\end{proposition}

The last assertion follows pointwise by deleting contributions.  We
will also use the following form, valid for every prescribed final
distance $0<e<\rho_{\rm col}$ with constants independent of $e$:
\begin{equation}
 \mathbb P\left\{\int_AW_e\,dA>u\right\}
       \le C\exp(-cu^{2/p}),\qquad u\ge u_0.
 \label{mesh:normalized-tail}
\end{equation}
Here is the exact passage from \eqref{mesh:tail-equation} to this
form.  Write $C_{\rm t}$ for the threshold constant in that equation.
The pointwise bound $W_e\le e^{-p}$ at a first live approach gives
$\int_AW_e\,dA\le C_{\rm d}e^{-p}$.  Choose
$C_0\ge\max(C_{\rm t},C_{\rm d})$, and for
$u\ge u_0:=C_0\delta_0^{-p}$ set
$\delta=(C_0/u)^{1/p}$.  If $e\le\delta$, then
\eqref{mesh:tail-equation} applies and its threshold is at most $u$.
If $e>\delta$, the deterministic bound is strictly less than
$C_{\rm d}\delta^{-p}\le u$, so the event is empty.  These two cases
prove \eqref{mesh:normalized-tail}.  Scaling space by $R$ multiplies
the density by $R^{-p}$ and area by $R^2$.  Thus on a disk of radius
$R$, the same bound has threshold $uR^d$ and permits final distances
$0<e<R\rho_{\rm col}$.

\subsection{Brownian occupation of a boundary layer}

The bracket estimate will use a geometric fact about the Dirichlet
Green function.  The Green occupation identity is classical; we use
the killed Brownian version described in
\cite[\S3.3, Definitions~3.28 and~3.31 and Theorem~3.32, pp.~78--80]{mortersperes2010}.  The uniform boundary-layer
estimate below requires an additional argument and no smoothness of
the boundary.

\begin{lemma}[Boundary-layer occupation]\label{mesh:layer}
Let $D$ be a simply connected plane domain with connected unbounded
complement, and let $K_D$ be its Dirichlet Green function in the
normalization
\[
 K_{\mathbb H}(z,w)=\log\left|\frac{z-\overline w}{z-w}\right|.
\]
For every $r>0$ and every $z\in D$,
\begin{equation}
 \int_{\{w\in D:\operatorname{dist}(w,\partial D)\le r\}}
 K_D(z,w)\,dA(w)\le Cr^2.
 \label{mesh:layer-bound}
\end{equation}
\end{lemma}

\begin{proof}
Figure~\ref{fig:brownian-circuit} illustrates the forced circuit and
the occupation blocks used below. We first record a uniform exit probability.  If $x\in D$ is within
distance $r$ of $D^c$, planar Brownian motion starting at $x$ has
probability at least $\eta>0$ of exiting $D$ before time $Ar^2$, for
absolute constants $A,\eta$.  To see this, scale and translate so that
$x=0,r=1$.  Brownian motion has a fixed positive probability of making,
within a fixed time, a circuit in a fixed annulus that separates
$\overline B(0,2)$ from infinity.  This is a finite open crossing event:
use overlapping rectangles along a circular route and two transverse
crossings of the final rectangle to force intersection with the
initial portion, and hence a separating closed subcurve.  Its
probability is positive by Brownian scaling, the Markov property and
positive transition probabilities through the prescribed rectangles.
If the path has not exited earlier, that circuit must meet $D^c$:
$D^c$ is connected, meets $\overline B(0,1)$, and is unbounded.  This
proves the asserted exit probability.

Let $\zeta$ be Brownian exit time from $D$, and let $E_r$ be the layer
in the integral.  Start a time block of length $Ar^2$ at the first
entrance into $E_r$.  After a block ends, start the next block at the
next entrance into $E_r$, unless exit has already occurred.  Entrance
times are stopping times; one may equally first use the open layer of
width $2r$ to avoid any convention at its outer boundary.  Conditional
on every block start before $\zeta$, the probability of exit during
the block is at least $\eta$.  Thus the expected number of blocks is
at most $1/\eta$.  All time spent in $E_r$ before exit is included in
these blocks, since the gaps contain no visit to $E_r$.  Consequently
\[
 \mathbb E_z\int_0^\zeta\mathbf1_{E_r}(B_t)\,dt
 \le Ar^2/\eta.
\]
The Green occupation formula gives \eqref{mesh:layer-bound}, up to the
absolute constant determined by the Brownian generator.  This formula
can first be applied in bounded exhaustions of $D$ and then passed to
the limit by monotone convergence.  The occupation estimate itself
shows that the layer integral remains finite.
\end{proof}

\begin{figure}[tbp]
\centering
\begingroup
\definecolor{circuitblue}{RGB}{31,88,138}
\definecolor{circuitorange}{RGB}{181,85,22}
\definecolor{circuitgray}{RGB}{104,110,116}
\begin{tikzpicture}[x=1cm,y=1cm,font=\small,
  line cap=round,line join=round,>=Latex]
\path[use as bounding box] (-.25,-.15) rectangle (14.6,10.1);
\node[anchor=west,font=\small\bfseries] at (0,9.85)
  {(a) A forced circuit after scaling $x=0$, $r=1$};
\node[anchor=west,font=\small\bfseries] at (8.25,9.85)
  {(b) The two crossings of $Q$};

\begin{scope}[shift={(3.35,6.2)},scale=.88]
  \fill[black!4] (0,0) circle (2);
  \draw[black!40,dashed] (0,0) circle (2);
  \draw[black!35,dotted] (0,0) circle (1);
  \fill (0,0) circle (1.6pt);
  \node[anchor=south east] at (-.08,.06) {$0$};
  \node at (-.15,1.2) {$\overline B(0,2)$};
  \draw[circuitblue!12,line width=12pt,rounded corners=12pt]
    (3,.6)--(3,2.4)--(2.4,3.15)--(-2.4,3.15)--(-3.15,2.4)
    --(-3.15,-2.4)--(-2.4,-3.15)--(2.4,-3.15)--(2.55,-2.25)
    --(2.55,0)--(3.55,0);
  \draw[circuitblue,line width=1.05pt]
    (3,.5)--(3.06,1.3)--(2.94,1.8)--(3,2.35)
    --(2.73,2.75)--(2.37,3.09)--(1.8,3.12)--(1.1,3.06)
    --(.35,3.2)--(-.25,3.08)--(-1,3.18)--(-1.72,3.11)
    --(-2.4,3.14)--(-2.82,2.76)--(-3.1,2.35)--(-3.18,1.65)
    --(-3.07,.8)--(-3.2,.12)--(-3.12,-.55)--(-3.16,-1.35)
    --(-3.1,-2.4)--(-2.72,-2.85)--(-2.36,-3.1)
    --(-1.6,-3.18)--(-.7,-3.08)--(.15,-3.15)--(.8,-3.1)
    --(1.5,-3.18)--(2.37,-3.1)--(2.58,-2.25)
    --(2.51,-1.5)--(2.6,-.8)--(2.55,0)--(2.62,0);
  \draw[circuitblue,-{Latex[length=2mm]}] (-.45,3.13)--(-.95,3.15);
  \draw[circuitblue,-{Latex[length=2mm]}] (.1,-3.14)--(.65,-3.12);
  \draw[circuitgray,dashed,thin]
    (0,0) .. controls (1,-.5) and (2.75,-1.25) .. (3,-.6);
  \node[anchor=north west,text=circuitgray,font=\footnotesize]
    at (-.25,-.68) {initial approach};
  \fill[white] (2.6,-.5) rectangle (3.4,.5);
  \draw[black!65,dashed] (2.6,-.5) rectangle (3.4,.5);
  \draw[circuitblue,line width=1.15pt,-{Latex[length=1.6mm]}]
    (3,-.65)--(2.94,-.3)--(3,0)--(3.08,.28)--(3,.66);
  \draw[circuitorange,line width=1.15pt,-{Latex[length=1.6mm]}]
    (2.49,.02)--(2.75,-.03)--(3,0)--(3.28,.07)--(3.55,.02);
  \node[anchor=west] at (3.45,.6) {$Q$};
  \node[anchor=north,align=center,font=\footnotesize]
    at (0,-3.5) {ordered corridor: one counterclockwise turn};
\end{scope}

\begin{scope}[shift={(11.2,6.6)}]
  \fill[circuitblue!6] (-.4,-1.65) rectangle (.4,1.65);
  \fill[circuitorange!8] (-1.65,-.4) rectangle (1.65,.4);
  \draw[black!70] (-1.65,-1.65) rectangle (1.65,1.65);
  \node[anchor=north east] at (1.62,1.64) {$Q$};
  \draw[circuitblue,line width=1.4pt,-{Latex[length=2mm]}]
    (-.05,-1.85)--(.13,-1.42)--(-.1,-.95)--(.15,-.6)
    --(-.1,-.27)--(0,0)--(.15,.37)--(-.16,.8)--(.08,1.25)--(0,1.9);
  \draw[circuitorange,line width=1.4pt,-{Latex[length=2mm]}]
    (-1.88,.06)--(-1.37,-.09)--(-.91,.14)--(-.46,-.12)
    --(0,0)--(.46,.12)--(.93,-.07)--(1.37,.09)--(1.89,.02);
  \fill (0,0) circle (2.1pt);
  \node[anchor=south west,align=left,font=\footnotesize]
    at (.36,.43) {a common point};
  \draw[thin] (.38,.44)--(.05,.06);
  \node[anchor=south,text=circuitblue] at (0,2.1) {first: bottom to top};
  \node[anchor=north,text=circuitorange,align=center]
    at (0,-1.98) {last: left to right};
\end{scope}

\node[anchor=west,font=\small\bfseries] at (0,2.28)
  {(c) Covering all occupation of $E_r$ by time blocks};
\draw[-{Latex[length=2mm]},black!70] (.2,.87)--(14.1,.87)
  node[anchor=west] {$t$};
\foreach \a/\b in {.8/3.5,5.0/7.7,9.2/11.9}{
  \fill[circuitblue!10] (\a,.91) rectangle (\b,1.43);
  \draw[circuitblue!55] (\a,.91) rectangle (\b,1.43);
}
\fill[white,opacity=.65] (10.95,.89) rectangle (12,1.48);
\draw[black!80,dashed] (10.95,.6)--(10.95,1.65);
\node[anchor=south] at (10.95,1.67) {$\zeta$};
\foreach \a/\b in {.8/1.35,1.8/2.0,2.65/3.0,5/5.85,6.7/7.3,9.2/9.75,10.1/10.8}
  \draw[circuitorange,line width=3pt] (\a,1.14)--(\b,1.14);
\draw[<->,thin] (.8,1.7)--(3.5,1.7)
  node[midway,fill=white,inner sep=1pt] {$Ar^2$};
\foreach \a/\lab in {.8/{\sigma_1},5/{\sigma_2},9.2/{\sigma_3}}
  \node[anchor=north] at (\a,.68) {$\lab$};
\node[anchor=south,align=center,font=\footnotesize] at (4.35,1.64)
  {outside $E_r$};
\draw[thin] (4.35,1.6)--(4.35,1.05);
\node[anchor=west,text=circuitorange,font=\footnotesize] at (11.7,.27)
  {visits to $E_r$};
\draw[circuitorange,line width=3pt] (11.15,.39)--(11.6,.39);
\node[anchor=west,font=\footnotesize,text=black!65] at (0,-.03)
  {Schematic; spatial corridors and time blocks are not to scale.};
\end{tikzpicture}
\endgroup

\caption{\BrownianCircuitCaption}
\label{fig:brownian-circuit}
\end{figure}
\clearpage

\subsection{Asynchronously started brackets for
\texorpdfstring{$\kappa\le4$}{kappa at most four}}

Schramm--Sheffield use Dirichlet Green decrease in their SLE$_4$
and Gaussian-free-field covariance calculation
\cite[Lemmas~2.5--2.6]{schrammsheffield2013}.
Here each point starts and stops at its own first approach times, so
the spatial bracket estimate needs the additional later-start
decomposition below.

Fix $0<q\le1/(8\mathrm e)$.  For the moment consider one
slab from first distance $r$ to first distance $qr$.  A point is
\emph{active} after $\tau_r(z)$ until $\tau_{qr}(z)$ or killing.  The
one-point martingale is bounded on this interval by $C r^{-p}$, where
$C$ may depend on the fixed ratio $q$.  Because $a_0\ge1$, its Brownian
coefficient $b_t(z)$, set to zero outside this interval, satisfies
\begin{equation}
 |b_t(z)|\le C r^{-p}\frac{Y_t(z)}{|Z_t(z)|^2}.
 \label{mesh:active-coefficient}
\end{equation}
Indeed the initial collar and the fact that the smaller trace
distance has not yet been reached give $R_{D_t}(z)\ge qr$, whereas
$\sin^{a_0}\theta\le\sin\theta$.

For two live distinct points put $K_t(z,w)=K_{D_t}(z,w)$.  The common
Brownian translation cancels in the differences appearing in its
half-plane formula.  Direct differentiation of the Loewner equation
therefore gives
\begin{equation}
 -\partial_tK_t(z,w)
 =4\frac{Y_t(z)}{|Z_t(z)|^2}\frac{Y_t(w)}{|Z_t(w)|^2}.
 \label{mesh:green-decrease}
\end{equation}
Consider pairs whose active intervals overlap.  Assign such a pair to
its later starting point, say $z$, so that
$\tau_r(w)\le\tau_r(z)$.  The contribution of the overlapping
interval to the integral of $|b_t(z)b_t(w)|$ is at most
\begin{equation}
 Cr^{-2p}K_{D_{\tau_r(z)}}(z,w).
 \label{mesh:pair-bound}
\end{equation}
Here we used positivity of the Green function at the end of the
overlap.  Infinite terminal intervals follow by monotone convergence.

At time $\tau_r(z)$ the earlier point $w$ is still in the domain and
lies within distance $r$ of its boundary.  This is because it was
within that distance at its own start and the domains decrease.
For each fixed $z$, its starting time does not depend on the dummy
integration variable $w$.  Lemma~\ref{mesh:layer} thus applies in the
single domain $D_{\tau_r(z)}$ and gives
\[
 \int_{\{w:\text{$w$ is active at $\tau_r(z)$}\}}
 K_{D_{\tau_r(z)}}(z,w)\,dA(w)\le Cr^2.
\]
Integrating in $z\in A$ and allowing both orders of the pair yields
the pathwise bound
\begin{equation}
 \int_0^\infty\left|\int_A b_t(z)\,dA(z)\right|^2dt
 \le\int_{A^2}\int_0^\infty |b_t(z)b_t(w)|\,dt\,dA(z)dA(w)
 \le Cr^{2-2p}.
 \label{mesh:bracket-bound}
\end{equation}
The spatial diagonal has area measure zero.  Ties of the starting
times may be assigned either order or counted twice; the upper bound
is unchanged.  Points with no start contribute zero.

We spell out the stochastic integration step.  For $s\in\{r,qr\}$,
let
\[
 M_t^{(s)}(z)=M_{t\wedge\tau_s(z)}(z),
 \qquad N_t(z)=M_t^{(qr)}(z)-M_t^{(r)}(z),
\]
using the zero continuation of $M$ at killing.  The collar and the
distance cutoff give $0\le M_t^{(s)}\le Cs^{-p}$.
If the distance approach never occurs, the live conformal radius
stays above $s$.  The radius terminal-zero assertion of
Lemma~\ref{ll:approach} then makes the zero continuation continuous
at killing and gives $M_\infty^{(s)}=0$.  At a finite approach the
process is frozen at $W_s$, so in all cases $M_\infty^{(s)}=W_s$.
Localization and bounded convergence therefore make each
stopped process a bounded continuous martingale.  Since
$\tau_{qr}\ge\tau_r$, the difference $N_t(z)$ is zero before its
start and has coefficient $b_t(z)$ precisely while the point is
active.  This formula also covers a missing start and an infinite
terminal interval.

Choose a version of $b_t(z)$ jointly measurable for the product of
the Borel sigma field on $A$ and the predictable sigma field.
Compact localization of the continuous adapted Loewner coefficients
and the indicators of the stochastic intervals supplies this version.  The
uniform bound $|N_t(z)|\le Cr^{-p}$ bounds its expected quadratic
variation by $Cr^{-2p}$.  On the finite area set $A$, this gives
\[
 \int_A\left(\mathbb E\int_0^\infty |b_t(z)|^2dt\right)^{1/2}
 dA(z)<\infty.
\]
The $L^2$ stochastic Fubini theorem, or simple spatial approximation
and the It\^o isometry, now identifies
$N_t=\int_A N_t(z)dA(z)$ with the Brownian integral of
$\int_A b_t(z)dA(z)$.  It converges in $L^2$ at infinity.  The same
uniform bound permits the pointwise terminal identity to pass through
the spatial integral, giving
\[
 N_\infty=\int_A(W_{qr}-W_r)\,dA.
\]
All uses of Fubini in \eqref{mesh:bracket-bound} concern nonnegative
functions and are valid pathwise.  Thus no interchange of
uncontrolled signed infinite-time integrals is needed.

The exponential martingale inequality and
\eqref{mesh:bracket-bound} imply
\begin{equation}
 \mathbb P\left\{\int_A(W_{qr}-W_r)\,dA>x\right\}
 \le\exp\left(-\frac{c x^2}{r^{2-2p}}\right).
 \label{mesh:slab-tail}
\end{equation}
This estimate, notably, permits different starting times for all
points in the spatial integral.

\subsection{Summing the slabs for
\texorpdfstring{$\kappa\le4$}{kappa at most four}}

Given $r_f\le\delta$, choose $N\ge0$ and
$r_0\in(q\delta,\delta]$ so that $r_j=r_0q^j$ and $r_N=r_f$.
This harmless adjustment ensures that even the final slab has ratio
exactly $q$.  The initial bound
$S_0:=\int_AW_{r_0}\,dA\le C\delta^{-p}$ is deterministic.
Apply \eqref{mesh:slab-tail} with
$x_j=\delta^{-p}/(j+1)^2$.  The exceptional probability is at most
\[
 \exp\left(-c\delta^{-2}
       \frac{q^{-2(1-p)j}}{(j+1)^4}\right).
\]
The sum of these probabilities over all $j\ge0$ is bounded by
$C\exp(-c'\delta^{-2})$: an exponential sequence divided by a fixed
power has a positive minimum and then grows geometrically.  Outside
this union of exceptions,
\[
 S_N\le S_0+\sum_{j=0}^{N-1}x_j\le C'\delta^{-p}.
\]
This proves Proposition~\ref{mesh:tail} for $\kappa\le4$.


\subsection{An integrated concentration estimate for disjoint windows}

The following elementary concentration argument allows the order of the
tests to depend on the curve.  In particular, it does not condition on a
permutation of the test locations.

\Needspace{26\baselineskip}
\begin{lemma}[Disjoint-window concentration]\label{win:concentration}
Let $D\subset B(0,1)$ be Borel, let $0<r\le R$ and $K>0$ with $3R\le1$, and let
$\gamma$ be a continuous adapted curve.  For each $z\in D$, suppose that
a test starts at a stopping time $\sigma_z$, with $\sigma_z=\infty$ if it
never starts.  At a finite start,
\[
 |\gamma(\sigma_z)-z|=r.
\]
Suppose its ending time $\tau_z\ge\sigma_z$ is a stopping time and
$\gamma(t)\in B(z,R)$ for $\sigma_z\le t<\tau_z$.
Let $Y_z$ be an $\mathcal F_{\tau_z}$-measurable output satisfying
$0\le Y_z\le K$, and let $\mu_z\ge0$ be $\mathcal F_{\sigma_z}$-measurable,
with
\[
 \mathbb E[Y_z\mid\mathcal F_{\sigma_z}]\le\mu_z
 \quad\hbox{on }\{\sigma_z<\infty\}.
\]
Assume the outputs and upper means are jointly measurable in $(z,\omega)$,
are zero when there is no start, and are uniformly bounded.  If a started
window never ends, assume its output is zero.  Then
\begin{equation}
 \mathbb E\exp\left(\lambda\int_D(Y_z-\mu_z)\,dA(z)\right)
       \le \exp(CK^2\lambda^2 R^2),\qquad \lambda\ge0,
 \label{win:mgf}
\end{equation}
and hence, for $b>0$,
\begin{equation}
 \mathbb P\left\{\int_D(Y_z-\mu_z)\,dA(z)>b\right\}
       \le \exp\left(-\frac{c b^2}{K^2R^2}\right).
 \label{win:tail}
\end{equation}
Here $C,c>0$ are absolute constants.  The same conclusion holds
conditionally on an initial sigma field if all the hypotheses and
conditional-mean estimates hold after that conditioning.
\end{lemma}

\begin{proof}
Set $a=3R$.  Fix a shift $\xi\in Q=[0,a)^2$, and consider the finitely
many centers of $(\xi+a\mathbb Z^2)\cap D$.  Their number is at most
$N\le C a^{-2}$.  Use the convention
$\mathcal F_\infty=\sigma(\bigcup_{t<\infty}\mathcal F_t)$.
While a window at $z$ is running, a window at any
other lattice center $w$ cannot start: otherwise
\[
 |z-w|\le |z-\gamma(\sigma_w)|+|\gamma(\sigma_w)-w|
          <R+r\le2R<3R.
\]
Order the finite starts by always waiting for the next start among the
remaining centers.  The minimum of the remaining stopping times is a
stopping time, and its index is measurable at that time.  Distinct finite
starts cannot tie, by the same geometric argument.  The preceding
window has ended by the next start, so its output is already known.
This is an adaptive enumeration, not conditioning on an eventual order.

For its $k$th actual test write $Z_k=Y_k-\mu_k$.  Conditional Hoeffding's
lemma gives, for $\theta\ge0$,
\[
 \mathbb E[e^{\theta Z_k}\mid\mathcal F_{\sigma_k}]
 \le \exp(K^2\theta^2/8).
\]
Indeed $Y_k$ has conditional range of length at most $K$, and
$\mathbb E[Y_k\mid\mathcal F_{\sigma_k}]-\mu_k\le0$.
For the selected test this inequality follows by partitioning over
the finitely many possible center indices, each known at the start.
The filtration at the end of test $k-1$ is contained in that at the
start of test $k$.  Iterating conditional expectations therefore yields
\[
 \mathbb E\exp\left(\theta\sum_{k=1}^{N}Z_k\right)
       \le \exp(NK^2\theta^2/8).
\]
Here missing tests are padded with zero increments.  This also covers
infinite waiting.  If a started window lasts forever, its confinement
prevents every later lattice start, its own increment is $-\mu_k\le0$,
and all later increments are zero.  Equivalently, one can use terminal
sigma fields at the infinite stopping times in the same argument.

Extend $f(z)=Y_z-\mu_z$ by zero outside $D$, and set
\[
 T_\xi=a^2\sum_{z\in\xi+a\mathbb Z^2}f(z).
\]
Taking $\theta=\lambda a^2$ gives
$\mathbb E e^{\lambda T_\xi}\le\exp(CK^2\lambda^2a^2)$.
The exact averaging identity is
\[
 \int_D f(z)\,dA(z)=\frac1{a^2}\int_Q T_\xi\,d\xi.
\]
Jensen's inequality, followed by Tonelli's theorem, proves
\eqref{win:mgf}.  There is no probabilistic independence assumption
on the shift.  Optimizing the exponential Markov bound in $\lambda$
proves \eqref{win:tail}.
\end{proof}

\subsection{The mesh tail for
\texorpdfstring{$4<\kappa<8$}{kappa between four and eight}}

\begin{proof}[Proof of Proposition~\ref{mesh:tail} for $4<\kappa<8$]
Use the first approach times $\tau_r(z)$, the densities $W_r(z)$,
and the indicators $I_r(z)$ already defined. The pointwise bound is
$W_r\le Cr^{-p}I_r$. Lemma~\ref{ll:approach} gives, at a finite
$\tau_r(z)$ and for one fixed $0<q\le1/(8\mathrm e)$,
\begin{equation}
 \mathbb P\{I_{qr}(z)=1\mid\mathcal F_{\tau_r(z)}\}
       \le C(qr)^pW_r(z).
 \label{win:onepoint-input}
\end{equation}
The conditional expected density at that finer approach is $W_r(z)$;
requiring it to occur before an additional stop can only decrease this
expectation. Proposition~\ref{ll:disk-rule} gives constants
$C_*,c_*,L_0>0$ for which a live exit from $B(z,Lr)$ has conditional
probability at most $C_*L^{-c_*}$ for $L\ge L_0$.
These are exactly the mean bounds needed for the three window outputs
below. Their constants are uniform over the permitted initial pasts.

Choose $r_0\in[q\delta,\delta]$ and an integer $n\ge0$ such that
$r_f=r_0q^n$, and set $r_j=r_0q^j$.  Thus every consecutive ratio is
exactly $q$, including the last; the starting mesh changes only by a
fixed factor.  Fix $0<\varepsilon<1-p$.  For the moment let
$L_*\ge L_0$ be arbitrary and set
\[
 L_j=L_*q^{-\varepsilon j}.
\]
During the estimates below suppose that $\delta<\rho_{\rm col}$ and
$3L_*\delta\le1$.  After deriving the recurrences we
will choose $L_*$ and then $\delta_0$ to enforce these requirements.
At every level,
\[
 L_jr_j=L_*r_0q^{(1-\varepsilon)j}\le L_*\delta\le\tfrac13,
\]
as required in Lemma~\ref{win:concentration}.  Also each
$\overline B(z,r_j)$ lies in the original collared initial domain, and a finite
$\tau_{r_j}(z)$ is its first circle hit with $z$ live.  These are the
hypotheses on the inner disk in Proposition~\ref{ll:disk-rule} at every level.

At a finite start $\tau_{r_j}(z)$, stop the window at the first of
the next live $r_{j+1}$-approach, exit from $B(z,L_jr_j)$, and death
of $z$.  Let $H_j(z)$ indicate a next approach achieved during this
window, let
\[
 V_j(z)=r_j^pW_{r_{j+1}}(z)H_j(z),
\]
and let $B_j(z)$ indicate an outer exit before that next approach,
with $z$ still live.  Set all three outputs to zero if no window
starts.  A window which never ends has all three outputs zero.
The outputs are bounded by a constant depending only on the fixed
parameters, since $V_j\le Cq^{-p}$.
First exits of the continuous process $t\mapsto|\gamma(t)-z|$ and death times
from the live Loewner maps are jointly measurable after compact
localization.  Thus these outputs and the upper means below have the
jointly measurable versions required by Lemma~\ref{win:concentration}.

The three conditional upper means at the start, also set to zero for
no start, are respectively
\begin{equation}
 \begin{split}
 \mu^H_j(z)&=C(qr_j)^pW_{r_j}(z),\\
 \mu^V_j(z)&=r_j^pW_{r_j}(z),\\
 \mu^B_j(z)&=C_*L_j^{-c_*}I_{r_j}(z).
 \end{split}
 \label{win:means}
\end{equation}
The pointwise bound on $W_{r_j}$ and $L_j\ge L_0$ make these upper
means uniformly bounded as well.
The first two follow from Lemma~\ref{ll:approach}.  The bad event is a
subset of the local-rule event, so the last bound follows from Proposition~\ref{ll:disk-rule}.  Lattice windows have disjoint time intervals by the
tip-containment argument in Lemma~\ref{win:concentration}, including
when a window ends early by killing.

For each of the three outputs apply that lemma at level
\[
 b_j=\frac{(r_j/\delta)^p}{(j+1)^2}\le1.
\]
Except on an event of probability at most
\begin{equation}
 C\exp\left(-\frac{c b_j^2}{(L_jr_j)^2}\right),
 \label{win:level-failure}
\end{equation}
the corresponding integral is at most the integral of its upper mean
plus $b_j$.  The exponent satisfies
\begin{equation}
 \frac{b_j^2}{(L_jr_j)^2}
 =\frac{\delta^{-2}}{L_*^2}
   \left(\frac{r_0}{\delta}\right)^{-2+2p}
   \frac{q^{-2(1-p-\varepsilon)j}}{(j+1)^4}
 \ge \frac{\delta^{-2}}{L_*^2}
   \frac{q^{-2(1-p-\varepsilon)j}}{(j+1)^4}.
 \label{win:failure-exponent}
\end{equation}
Because $q^{-2(1-p-\varepsilon)}>1$, the sum of
\eqref{win:level-failure} over all $j\ge0$ and the three outputs is
at most $C\exp(-c\delta^{-2}/L_*^2)$.  For example, the last quotient in
\eqref{win:failure-exponent} is bounded below by a positive constant
times $j+1$, which reduces the sum to a geometric series.  This bound
is independent of $n$ and hence of $r_f$.

It remains to propagate the simultaneous integral bounds.  A next
approach missed by its window must occur after an outer exit at which
the center was still live.  Death or an infinite window cannot be
followed by a missed live approach.  Consequently, pointwise,
\begin{equation}
 \begin{split}
 I_{r_{j+1}}&\le H_j+B_j,\\
 W_{r_{j+1}}&\le r_j^{-p}V_j+C(qr_j)^{-p}B_j.
 \end{split}
 \label{win:missed-output}
\end{equation}
Define
\[
 S_j=\int_A W_{r_j}\,dA,\qquad
 A_j=r_j^{-p}\int_A I_{r_j}\,dA.
\]
On the simultaneous good event, integrating
\eqref{win:missed-output} and using \eqref{win:means} gives
\begin{equation}
 \begin{split}
 S_{j+1}&\le S_j+\alpha_j A_j+e_j,\\
 A_{j+1}&\le C_1 S_j+\alpha_j A_j+e_j,
 \end{split}
 \label{win:recurrences}
\end{equation}
where constants can be enlarged so that
\[
 \alpha_j=C_2L_*^{-c_*}q^{\varepsilon c_*j},\qquad
 e_j=\frac{C_3\delta^{-p}}{(j+1)^2}.
\]
Here $r_j^{-p}b_j=\delta^{-p}/(j+1)^2$ exactly; factors of $q^{-p}$
are fixed constants.  The constants $C_1,C_2,C_3$ depend only on
$q$ and the constants in the local laws, not on $L_*$ or the final mesh.
The pointwise initial bounds and
$r_0\ge q\delta$ give $S_0,A_0\le C\delta^{-p}$.

For completeness, the two recurrences have a bound uniform in their
length.  Choose $\Lambda\ge\max(1,2C_1)$ and then $L_*$ so large that
$\alpha_0=C_2L_*^{-c_*}\le1/2$.  Finally choose
$0<\delta_0<\rho_{\rm col}$ with $3L_*\delta_0\le1$.  The preceding
estimates now apply whenever $r_f\le\delta\le\delta_0$.
Put $U_j=\max(S_j,A_j/\Lambda)$.  The second recurrence
implies
\[
 A_{j+1}/\Lambda
 \le(C_1/\Lambda+\alpha_j)U_j+e_j/\Lambda\le U_j+e_j,
\]
while the first implies
$S_{j+1}\le(1+\Lambda\alpha_j)U_j+e_j$.  Thus
\[
 U_{j+1}\le(1+\Lambda\alpha_j)U_j+e_j,
\]
and iteration gives
\[
 U_j\le\left(U_0+\sum_{i\ge0}e_i\right)
          \exp\left(\Lambda\sum_{i\ge0}\alpha_i\right)
       \le C\delta^{-p}.
\]
Both sums converge and their constants are independent of the final
mesh.  With $L_*$ fixed, the summed failure probability is
$C\exp(-c\delta^{-2})$.  Thus $S_n\le C\delta^{-p}$ outside that event,
proving \eqref{mesh:tail-equation} for $4<\kappa<8$ and completing the proposition.
\end{proof}
